%=====================================================================
\section{Family generator: \texttt{ctvae} (contract 1.1)}\label{sec:v74ctvae}
% numbered from 1860; the shared counter is restored at the end
\newcounter{savedpropctvae}\setcounter{savedpropctvae}{\value{proposition}}
A \emph{family} is the set of rows of a fact that share an anchor row (the lines of an invoice, the transactions of an account). Row-wise generators draw a family's rows independently given each row's context. The generator \texttt{ctvae} declares the \texttt{family} conditioning of contract 1.1 and draws the rows of a family jointly. It has three parts: a row VAE with a PreNorm transformer over column tokens; a causal transformer (or GRU) prior over the row latents, with a mixture density network for $p(z_i\mid\text{history})$; and a residual categorical head for $p(c_i\mid\text{history},z_i)$. The columns are decoded through the Bayesian marginals, so Props.~1800--1805 apply unchanged. The design and the ablation that chose it (four datasets, three seeds, acceptance rule fixed before the runs) are in \texttt{docs/research/context-transformer-vae.md} and \texttt{experiments/context\_vae/}. The code is \texttt{relsynth/generators/neural/ctvae.py}, and the tests are in \texttt{tests/unit/test\_ctvae.py}.

\setcounter{proposition}{1859}
\begin{proposition}[The family law]\label{prop:ctvaefam}
Let a family have $n$ rows with contexts $a_1,\dots,a_n$ (segment, coarse segment, parent attributes), gaps $\gamma_1,\dots,\gamma_n$ ($\gamma_1=0$; all $0$ for unordered families) and rows $x_i=(y_i,c_i)$, where $y_i$ are the numeric scores and $c_i$ the categorical slots. The sampler draws, for $i=1,\dots,n$, $z_i\sim p_\theta(z\mid h_i)$, then $c_i\sim p_\theta(c\mid h_i,z_i)$ column by column, then $y_i\sim N(\mu_\psi(z_i),\mathrm{diag}\,s^2)$. Here $h_i=f_\theta(z_{<i},c_{<i},\gamma_{\le i},a_{\le i},\log n)$, restricted to the last $W$ rows. (a) This defines a proper joint law $p(x_{1:n},z_{1:n}\mid a,\gamma,n)=\prod_i p(z_i\mid h_i)\,p(c_i\mid h_i,z_i)\,p(y_i\mid z_i)$, and the rows of different families are independent given their contexts. (b) For unordered families trained on uniformly random permutations, the population minimiser of the expected teacher-forced log loss over all functions $f$ of the retained history $r_i$ (the last $W$ rows and the inputs above) is, for each $r_i$, the conditional law of the next row given $r_i$ under the mixture over orders. If the retained history is sufficient, i.e.\ the next row is conditionally independent of the dropped rows given $r_i$ (for instance $W\ge n$), these conditionals are the one-step conditionals of the order-averaged law of the pairs, $\frac1{n!}\sum_\pi p\big((a,x)_{\pi(1:n)}\big)$, read in a fixed order. When the fitted conditionals agree across orders, the sampled family is invariant in law under permuting the rows \emph{together with their contexts}; the attributes alone are exchangeable only when the contexts are identical (with $x_i=a_i$ and contexts $a=(0,1)$, the order-invariant predictor gives $(0,1)$ with probability one and $(1,0)$ with probability zero). Without sufficiency the minimiser is only the best predictor from $r_i$, and the joint law it generates need not be the order-averaged one. The gap between two orders on held-out families measures disagreement in either case (\texttt{experiments/context\_vae}, \texttt{order\_gap}).
\end{proposition}
\begin{proof}
(a) Each factor is a normalised density or mass function given the variables before it in the stated order (an MDN is a mixture of normals, a softmax a probability vector, the decoder noise Gaussian), so the product is the chain-rule factorisation of a joint law. The sampler draws the factors in that order. Families share parameters but no random variables, so they are conditionally independent. (b) Under random permutations the training pairs (history, next row) are drawn from the mixture over orders. The expected log loss splits over retained histories and is minimised by the conditional law of the next row given $r_i$ under that mixture (strict propriety of the log score). Under sufficiency the conditional given $r_i$ equals the conditional given the full history, so these are the one-step conditionals of the order-averaged law read in the order presented, and that law is invariant to permutations of the (context, row) pairs, which is all that averaging over orders gives. A finite window can merge histories with different next-row laws (with $W=1$, families uniform on $011$, $101$, $110$ have histories $01$ and $11$ with the same last row but next rows $1$ and $0$), which is why sufficiency is assumed.
\end{proof}
\noindent Ordered families (a fact with a time column) use the engine's dates as the order and the gaps as inputs; the engine draws a family's dates before its attributes (Prop.~1862).

\setcounter{proposition}{1860}
\begin{proposition}[The residual categorical head nests the decoder]\label{prop:ctvaeres}
Let the logits of categorical column $j$ be $\ell_j=d_j(z)+r_j(h+Uz+\textstyle\sum_{k<j}E_k[c_k])$, where $d_j$ are the row decoder's logits and $r_j$ is linear with weights and bias initialised at $0$. (a) At initialisation, $p(c_j\mid h,z,c_{<j})=\mathrm{softmax}(d_j(z))$, so the categorical part is the latent-only model's. (b) The family of conditionals contains the latent-only model ($r_j\equiv0$), so the minimum of the teacher-forced log loss over the parameters of $r_j$ (decoder fixed, stage 2) is at most the latent-only model's loss on the same data. (c) The correction can express dependence on earlier siblings' levels that is not carried by $z_i$: for instance, a level shared by siblings independently of the row's numerics.
\end{proposition}
\begin{proof}
(a) $r_j(\cdot)=0$ for zero weights and bias. (b) The parameter set of $r_j$ includes $0$, where the loss equals the latent-only loss; a minimum over a set that contains a point is no larger than the value at that point. (c) $h_i$ is a function of $c_{<i}$ through the input embeddings, and $r_j$ maps it to logits; with $U=0$ any function of $c_{<i}$ representable through $h$ can shift the logits.
\end{proof}
\noindent In the ablation, a latent-only prior lost categorical sharing on two of four datasets, and a head that replaced the decoder (no residual) lost it on the smallest dataset. The residual head kept or improved it on three of four (design record, run 4). Part (b) concerns the training loss; the held-out effect is what the ablation measured.

\setcounter{proposition}{1861}
\begin{proposition}[The engine's skeleton and other generators are unchanged]\label{prop:ctvaeskel}
(a) For a family generator the engine passes, at fit, the anchor row of every fitted fact row and its time column as order key, and at sampling the synthetic anchor row and synthetic dates. Keys, links, counts, dates, rules and uniqueness are set and post-processed by the engine as for any generator (Prop.~700). (b) Drawing the dates before the attributes leaves the law of the dates unchanged, because the engine's date law depends only on the anchor segment and not on the attributes; the attributes are then drawn from the generator's conditional law given the dates (for a generator that does not read the dates, the attribute law of contract 1.0). (c) A generator that does not declare \texttt{family} receives no family, and the engine's random streams are consumed in the same order as before contract 1.1, so its output for a seed is byte-identical.
\end{proposition}
\begin{proof}
(a) The generator returns only the table's attribute columns; the engine adds keys, dates and the post-processing afterwards, as for row-wise generators (\texttt{Synthesizer.\_fact\_rows}). (b) The dates are $t=\tau(z_1[u],\omega)$ for an independent randomness $\omega$, so their law is that of contract 1.0 whichever order the draws are made in, because $\tau$ does not read $x$. The attributes are $x=g(\text{context},t,\omega')$, so the released pair has the law of $t$ followed by $x\mid t$ under $g$. This does not claim that $x\mid t$ equals the contract 1.0 attribute law for a generator that reads $t$; it equals it when $g$ does not read $t$. (c) The reordering is taken only when \texttt{TableModel.uses\_family} is true. Otherwise the call sequence is the one of contract 1.0: attributes, then dates, from the same generator.
\end{proof}
\noindent Rows that the engine rejects (attribute rules, unique keys, near copies, row validity) are redrawn with their family ids, given the family's kept rows: the kept rows' latents and levels are fed to the prior in their order (teacher forcing) and only the redrawn rows are drawn, so the family size and history are the full family's. Test \texttt{test\_keys\_links\_and\_rules\_hold}; the byte identity of the other generators is checked by the bench's identity rule.

\setcounter{proposition}{1862}
\begin{proposition}[Entity-level DP-SGD]\label{prop:ctvaedp}
Fix the preprocessing (encoding marginals, kept levels, segment indices). Let $D$ be the set of encoded families, and let $N$ be a public unit count, the same for neighbouring inputs: \texttt{dp.population} when the spec sets it, else $|D|$, which is then treated as public (the usual assumption of DP-SGD, Abadi et al.\ 2016). The rate $q=\min(1,B/N)$, $\delta$, $\sigma$ and the gradient divisor $qN$ are functions of $N$ and the options only. Each of $T=T_1+T_2$ steps (row VAE, then prior) includes every family independently with probability $q$. It computes each included family's gradient over at most $W$ of its rows (chosen with randomness independent of $D$'s values), clips it to L2 norm $C$, sums, adds $N(0,\sigma^2C^2I)$ and takes an optimiser step that depends only on the noisy sums and earlier iterates; stage 2's inputs (posterior means) are functions of the stage-1 weights and the family's own rows. Then the weights, and every synthetic dataset sampled from them without the family calibration (the default under \texttt{dp}), are $(\varepsilon,\delta)$-DP with respect to adding or removing one whole family (one anchor row's rows in this table), with $\varepsilon$ the RDP bound of Prop.~1802 for $T$ steps at rate $q$ and noise multiplier $\sigma$.
\end{proposition}
\begin{proof}
In each step the clipped sum changes by at most one clipped family gradient when one family is added or removed, so the step is the sampled Gaussian mechanism with sensitivity $C$ on the unit ``family'' (Prop.~1802's argument with families in place of rows). A family's stage-2 gradient depends on $D$ only through its own rows and the weights released by earlier steps, which adaptive composition allows. Because $q$, $\sigma$ and the divisor $qN$ are functions of the public $N$ and not of $D$, neighbouring inputs run the same mechanism; the divisor is a fixed rescaling of the noisy sum, i.e.\ post-processing. The two stages are $T$ adaptively composed steps with the same $q$ and $\sigma$, so their R\'enyi divergences add (Mironov 2017, Prop.~1). The conversion to $(\varepsilon,\delta)$ and the post-processing are as in Prop.~1802: sampling reads only the weights, the fixed preprocessing and the synthetic skeleton the engine supplies.
\end{proof}
\noindent The unit is stronger than a row: one anchor row's whole history in the table is protected together. As for \texttt{tvae}, the bound is conditional on preprocessing fitted without noise and covers this table only. Which rows form a family comes from the engine's relationships. The family calibration (Prop.~1805 on families) draws from the weights for a random subset of the training families' sizes and contexts, which depends on whether one family is present; it is therefore off under \texttt{dp} unless the spec sets \texttt{calibrate: true}, and then it joins the preprocessing outside the bound. Without \texttt{dp.population} the family count itself is not protected: two inputs that differ by one family run mechanisms with slightly different $q$ (for $B=8$, $q=1$ at eight families and $8/9$ at nine), so the bound holds only if the count is public; setting \texttt{dp.population} to a public figure (for example a published number of accounts) removes the assumption. The generator reports \texttt{dp: false} with the spent $\varepsilon,\delta$; DP mode remains the end-to-end route. Test \texttt{test\_entity\_level\_dp\_sgd}.

\setcounter{proposition}{\value{savedpropctvae}}
